% Fourier Modal Solver: Architecture, Algorithms, and Benchmarks
% Source: docs/modal_solver.md (migrated to LaTeX)
% Uses macros from preamble.tex

\section{Fourier Modal Solver}
\label{sec:modal-solver}

\subsection{Overview}
\label{sec:modal-overview}

The Fourier modal solver (\texttt{tidal/solver/modal.py}) provides exact, machine-precision solutions for linear PDE systems with periodic boundary conditions and time-independent coefficients. It transforms the spatial grid to Fourier space, builds a per-mode evolution matrix, and computes $\exp(A \cdot t) \cdot y_0$ via Pad\'{e} matrix-exponential (\cref{sec:modal-robust-matrix-exp}) --- eliminating spatial discretization error entirely.

\textbf{Key advantages over PDE solvers (CVODE/IDA):}
\begin{itemize}
  \item Machine-precision accuracy (${\sim}10^{-14}$ error, limited only by IEEE-754 floating-point precision; see \cref{sec:modal-snapshot-precision})
  \item No CFL condition or adaptive timestepping needed
  \item $O(N^2)$ per output time after $O(N^3)$ precomputation
  \item Exact for any time point --- no time-stepping error accumulation
\end{itemize}

\textbf{Limitations:}
\begin{itemize}
  \item Requires periodic boundary conditions (FFT periodicity)
  \item Requires time-independent coefficients (evolution matrix must be constant)
  \item Requires flat (Minkowski) metric --- curved metrics couple modes via position-dependent coefficients that decay too slowly in $k$-space
  \item Default backend is pure NumPy/SciPy; a JAX-accelerated path
    (\texttt{--scheme modal-jax}) is available for PolyChord inference runs
    --- see \S\ref{sec:modal_jax}
\end{itemize}

\subsection{How It Works}
\label{sec:modal-how}

\subsubsection{First-Order Reformulation}
\label{sec:modal-first-order}

TIDAL stores state in Euler--Lagrange velocity form $(q, v)$ where $v = \mathrm{d}q/\mathrm{d}t$. A second-order PDE $\ddt^2 \varphi = L[\varphi]$ becomes:
\begin{align}
  \frac{\mathrm{d}\varphi}{\mathrm{d}t} &= v \\
  \frac{\mathrm{d}v}{\mathrm{d}t} &= L[\varphi]
\end{align}

The \texttt{StateLayout} organizes field+velocity slot pairs. The modal solver reads this directly.

\subsubsection{FFT to Fourier Space}
\label{sec:modal-fft}

Each field and velocity slot is transformed: $u(x) \to \hat{u}(k)$ via \texttt{numpy.fft.rfftn}. Spatial operators become multiplications:

\begin{table}[htbp]
\centering
\caption{Fourier multipliers for spatial operators.}
\label{tab:fourier-multipliers}
\begin{tabular}{@{}ll@{}}
\toprule
Operator & Fourier multiplier \\
\midrule
identity & $1$ \\
laplacian & $-k^2$ \\
laplacian\_x & $-k_x^2$ \\
gradient\_x & $ik_x$ \\
cross\_derivative\_xy & $-k_x k_y$ \\
biharmonic & $k^4$ \\
\bottomrule
\end{tabular}
\end{table}

\textbf{Fourier convention:} $k = 2\pi \cdot \texttt{numpy.fft.rfftfreq}(N, d=\Delta x)$ --- consistent with \texttt{operators.py} spectral mode. \emph{Not} the modified-wavenumber convention from \texttt{constraint\_solve.py}.

\subsubsection{Build Evolution Matrix}
\label{sec:modal-evolution-matrix}

\paragraph{Constant coefficients (all terms position-independent).}
\label{par:modal-constant}
Each $k$-mode decouples. For a 2-field second-order system, each mode yields a $4\times 4$ matrix:
\begin{equation}
  \frac{\mathrm{d}}{\mathrm{d}t}
  \begin{pmatrix} \hat\varphi_k \\ \hat v_{\varphi,k} \\ \hat\chi_k \\ \hat v_{\chi,k} \end{pmatrix}
  = M(k) \cdot
  \begin{pmatrix} \hat\varphi_k \\ \hat v_{\varphi,k} \\ \hat\chi_k \\ \hat v_{\chi,k} \end{pmatrix}
  \label{eq:modal-per-mode}
\end{equation}

Result: $N_\text{modes}$ independent small matrices. Total: $O(N_\text{modes} \times N_\text{fields}^3)$.

\paragraph{Position-dependent coefficients (convolution coupling).}
\label{par:modal-posdep}
When a coefficient $c(x)$ multiplies a field, the product in $k$-space becomes a convolution:
\begin{equation}
  \text{FFT}[c(x) \cdot u(x)] = \hat{c} * \hat{u} \quad\text{($k$-space convolution)}
  \label{eq:modal-convolution}
\end{equation}

This couples different $k$-modes, creating a full $(N_\text{modes} \times N_\text{slots})^2$ matrix. The convolution paths operate in the \textbf{full complex FFT basis} (all $N$ frequencies per axis, \texttt{numpy.fft.fftn}), where pointwise multiplication by $c(x)$ is exactly $\mathbb{C}$-linear and the block is the circulant $\hat{c}[(k-k') \bmod N]/N_\text{tot}$, built from a single FFT of $c$ (GH \#445). The earlier probe-vector construction over the rfft half-spectrum was retired there: the DC/Nyquist bins carry no imaginary degree of freedom in that basis, so sine-phase content that $c$'s harmonics fold onto them was mismapped (7--25\% per-mode action errors). Position-dependent \emph{kinetic} coefficients fold $M^{-1}(x)$ into the real-space coefficient before the FFT on the same paths (GH \#427). The constant-coefficient per-mode paths keep the cheaper rfft basis --- a constant coefficient is diagonal in $k$, so no folding occurs. For Gaussian $c(x)$, the kernel $\hat{c}(q)$ decays exponentially $\to$ banded matrix.

\subsubsection{Time Evolution Algorithms}
\label{sec:modal-time-evolution}

Two paths depending on coefficient structure:

\paragraph{Per-mode eigendecomposition (constant coefficients).}
\label{par:modal-eigendecomp}
When all coefficients are constant, the evolution matrix is block-diagonal in mode space. Each mode's small block ($\leq 12\times 12$) is eigendecomposed independently:
\begin{lstlisting}[style=python]
eigenvalues, V = np.linalg.eig(A_k)  # per-mode block
y_k(t) = V @ (exp(eigenvalues * t) * V_inv @ y0_k)
\end{lstlisting}

This has no time-stepping error and no CFL condition. $O(N_\text{fields}^3)$ per mode, $O(N_\text{modes} \times N_\text{fields}^3)$ total.

\paragraph{Nyquist mode zeroing (critical for energy conservation).}
\label{par:modal-nyquist}
The per-mode eigendecomposition achieves \textbf{machine-precision energy conservation} ($|\mathrm{d}E/E| \approx 2\times 10^{-14}$) after zeroing the Nyquist mode in the initial conditions. Without this, conservation degrades to ${\sim}1.5\times 10^{-5}$.

\textbf{Root cause}: The rfft Nyquist bin ($k = N/2$, the last mode in even-$N$ grids) must have a purely real Fourier coefficient for real-valued fields. However, the per-mode evolution matrix $A_k$ has complex entries from the gradient coupling ($ik$ term). When $\exp(A_k \cdot t)$ evolves the Nyquist mode, it creates imaginary components in the Fourier coefficient. The \texttt{irfft} reconstruction silently drops these imaginary parts (by construction --- the Nyquist bin of a real signal \emph{is} real). This truncation discards energy from the Nyquist mode at every snapshot.

\textbf{Investigation}: Per-mode analysis showed that \emph{all} modes except the Nyquist are exactly conserved to machine precision ($\mathrm{d}H_k = 0$ for $k < N/2$). The Nyquist mode's Hamiltonian grew by $500\times$ (from 0.029 to 16.8) in the complex eigendecomposition, but the stored (irfft-reconstructed) field showed shrinking amplitude --- the discarded imaginary part contained the ``missing'' energy.

\textbf{Fix}: Zero the Nyquist mode's IC before eigendecomposition. This is standard practice in pseudospectral methods --- the Nyquist frequency aliases with its negative-frequency conjugate and cannot faithfully represent directional (travelling wave) content. It is the highest resolvable frequency and carries negligible physical content for well-resolved simulations.

\textbf{Results} (coupled\_scalars, Gertsenshtein gradient coupling):

\begin{table}[htbp]
\centering
\caption{Energy conservation before and after Nyquist mode zeroing.}
\label{tab:nyquist-fix}
\begin{tabular}{@{}lcc@{}}
\toprule
Grid & Before fix & After fix \\
\midrule
$N=128$ & $|\mathrm{d}E/E| = 2.8\times 10^{-5}$ & $|\mathrm{d}E/E| = 2.2\times 10^{-14}$ \\
$N=256$ & $|\mathrm{d}E/E| = 1.8\times 10^{-5}$ & $|\mathrm{d}E/E| = 2.2\times 10^{-14}$ \\
$N=512$ & $|\mathrm{d}E/E| = 1.7\times 10^{-5}$ & $|\mathrm{d}E/E| = 2.2\times 10^{-14}$ \\
$N=1024$ & $|\mathrm{d}E/E| = 1.6\times 10^{-5}$ & $|\mathrm{d}E/E| = 2.2\times 10^{-14}$ \\
\bottomrule
\end{tabular}
\end{table}

Full Gertsenshtein (6 fields): $|\mathrm{d}E/E| = 1.3\times 10^{-14}$.

\textbf{References}:
\begin{itemize}
  \item Boyd, J.P.\ (2001). \emph{Chebyshev and Fourier Spectral Methods}, 2nd ed.\ Dover. \S11.5 (aliasing and the Nyquist mode).
  \item Canuto, C.\ et al.\ (2006). \emph{Spectral Methods: Fundamentals in Single Domains}. Springer. \S3.2 (dealiasing via mode truncation).
  \item Burns et al.\ (2020). Dedalus framework zeroes the top $1/3$ of modes (``$2/3$ dealiasing rule'') to prevent aliasing in nonlinear terms. Our Nyquist zeroing is the minimal version for linear systems.
\end{itemize}

\paragraph{Krylov matrix exponential (position-dependent coefficients).}
\label{par:modal-krylov}
Position-dependent convolution produces a full $(N_\text{slots} \times N_\text{modes})^2$ matrix that is generally \textbf{non-normal}: gradient operators ($ik$) combined with real convolution kernels create eigenvalues with significant positive real parts despite conservative physics. Eigendecomposition overflows because individual $\exp(\lambda \cdot t)$ diverge even though $\exp(A \cdot t) \cdot y_0$ is bounded (pseudospectral phenomenon; Trefethen \& Embree 2005, Ch.\ 14).

The fix: \texttt{scipy.sparse.linalg.expm\_multiply} computes $\exp(A \cdot t) \cdot y_0$ directly via scaling + truncated Taylor series in matrix-vector products, which is backward-stable for non-normal matrices:

\begin{lstlisting}[style=python]
from scipy.sparse.linalg import expm_multiply
y_all = expm_multiply(A_full, y0_flat, start=0,
                      stop=t_end, num=n_snapshots)
\end{lstlisting}

Ref: Al-Mohy \& Higham (2011), ``Computing the Action of the Matrix Exponential'', SIAM J.\ Sci.\ Comput.\ 33(2):488--511.

\subsubsection{Algorithm Selection (Dual-Path Design)}
\label{sec:modal-algorithm-selection}

The modal solver maintains two distinct time-evolution algorithms because they handle fundamentally different matrix structures. The routing is automatic --- \texttt{\_has\_position\_dependent\_terms()} determines which path is used at runtime. When a system has algebraic constraints or time-derivative operators on the RHS, a third unified reduction step (\texttt{\_build\_evolution\_matrices}) runs upstream of both paths to produce a reduced first-order evolution matrix.

\paragraph{Why two methods?}
\label{par:modal-why-two}

\textbf{Constant coefficients} (uniform $\Bzero$, mass terms, coupling constants): each Fourier mode decouples, producing small per-mode blocks (e.g.\ $12\times 12$ for 6-field Gertsenshtein). Eigendecomposition of these tiny matrices is exact, fast, and well-conditioned.

\textbf{Position-dependent coefficients} (localized $\Bzero(x)$, spatially-varying backgrounds): the product $c(x) \cdot u(x)$ becomes a convolution $\hat{c} * \hat{u}$ in $k$-space, coupling all modes into a single large matrix (e.g.\ $3000\times 3000$ for $N=512$, 6 fields). This matrix is \textbf{non-normal} --- gradient operators ($ik$) combined with real convolution kernels create eigenvalues with significant positive real parts despite conservative physics. Individual $\exp(\lambda \cdot t)$ diverge even though $\exp(A \cdot t) \cdot y_0$ is bounded (pseudospectral phenomenon; Trefethen \& Embree 2005, Ch.\ 14).

\textbf{Concrete failure case:} For localized Gertsenshtein, eigendecomposition gave $P = 0.477$ (wrong) vs the correct $P = 0.3437$ from the Boccaletti formula. The eigendecomposition was also slower (44s) than \texttt{expm\_multiply} (21s) for this case.

\paragraph{Algorithm routing.}
\label{par:modal-routing}

\begin{table}[htbp]
\centering
\caption{Algorithm routing for the modal solver's dual-path design.}
\label{tab:modal-routing}
\begin{tabular}{@{}lllll@{}}
\toprule
Case & Coefficient type & Matrix structure & Algorithm & Time (Gertsenshtein, $N=512$) \\
\midrule
Constant & Uniform & Per-mode blocks ($12\times 12$) & Eigendecomposition & ${\sim}0.02$\,s \\
Position-dep. & Spatially varying & Sparse convolution ($3084\times 3084$) & \texttt{expm\_multiply} (sparse CSC) & ${\sim}1.4$\,s \\
Constraints + const. & Uniform with algebraic constraints \emph{and/or} rank-deficient $M$ & Reduced per-mode blocks via unified Schur & \texttt{scipy.linalg.eig(A, B)} (QZ) & ${\sim}0.08$\,s \\
\bottomrule
\end{tabular}
\end{table}

Detection: \texttt{\_has\_position\_dependent\_terms()} checks for spatially-varying coefficients in the JSON spec. \textbf{Both paths are $O(1)$ in simulation time $t_\text{end}$} --- the key advantage over time-stepping solvers.

\paragraph{Can't we just use one method for everything?}
\label{par:modal-one-method}

\begin{itemize}
  \item \textbf{\texttt{expm\_multiply} for everything?} No --- for the small per-mode blocks ($bs \le 15$) of the constant-coefficient path, \texttt{expm\_multiply} carries per-call Krylov-Taylor setup overhead that vastly dominates the $O(bs^3)$ exponential cost. Precomputing $\exp(M \cdot \mathrm{d}t)$ once per mode via \texttt{scipy.linalg.expm} (Pad\'{e} path~d; see \cref{sec:modal-robust-matrix-exp}) is $0.49$--$1.06\times$ the wall time of the former eigendecomposition path while giving equivalent machine precision.
  \item \textbf{Eigendecomposition for everything?} No --- it gives \textbf{wrong results} for non-normal convolution matrices (position-dependent case). This is not a precision issue but a fundamental numerical instability: individual $\exp(\lambda \cdot t)$ overflow while $\exp(A \cdot t) \cdot y_0$ is bounded (pseudospectral phenomenon). The eigendecomposition path for the constant-coefficient (per-mode) case was also retired in v0.31 in favour of the unconditional Pad\'{e} path (\cref{sec:modal-pade-pass1}), which is robust for any $\mathrm{cond}(V)$.
\end{itemize}

\paragraph{Sparse convolution matrix optimization.}
\label{par:modal-sparse}

For localized coefficients (e.g.\ Gaussian $\Bzero(x)$), the Fourier convolution kernel $\hat{c}(q)$ decays exponentially, making the matrix effectively banded. The \texttt{\_evolve\_full\_matrix} function exploits this: entries below a relative threshold ($10^{-14} \times \max|A|$) are zeroed, and if the resulting density is below 30\%, the matrix is converted to sparse CSC format. \texttt{scipy.sparse.linalg.expm\_multiply} natively supports sparse matrices, and its internal matrix-vector products become dramatically cheaper.

\textbf{Benchmark: localized Gertsenshtein (6 fields, $t_\text{end}=120$, 51 snapshots)}

\begin{table}[htbp]
\centering
\caption{Sparse vs.\ dense \texttt{expm\_multiply} for localized Gertsenshtein.}
\label{tab:sparse-benchmark}
\begin{tabular}{@{}lcccccl@{}}
\toprule
Grid $N$ & Matrix size & Density & Dense & Sparse & Speedup & Max diff \\
\midrule
128 & $780\times 780$ & 4.0\% & 5.44\,s & 0.23\,s & $\mathbf{23\times}$ & $7.4\times 10^{-14}$ \\
256 & $1548\times 1548$ & 2.6\% & 15.85\,s & 0.59\,s & $\mathbf{27\times}$ & $1.2\times 10^{-13}$ \\
512 & $3084\times 3084$ & 1.4\% & 29.08\,s & 1.23\,s & $\mathbf{24\times}$ & $1.1\times 10^{-12}$ \\
\bottomrule
\end{tabular}
\end{table}

The density decreases with $N$ because the convolution bandwidth (set by the Gaussian width $R$) is fixed while the total mode count grows --- the matrix becomes sparser at higher resolution.

\textbf{End-to-end improvement:} Total \texttt{solve\_modal} time for localized Gertsenshtein at $N=512$ dropped from ${\sim}29$\,s (dense) to ${\sim}1.4$\,s (sparse + micro-optimizations), bringing it within ${\sim}70\times$ of the constant-coefficient eigendecomposition path and $\mathbf{{\sim}85\times}$ \textbf{faster than CVODE} (${\sim}120$\,s).

\paragraph{Both paths outperform CVODE.}
\label{par:modal-vs-cvode}

\begin{table}[htbp]
\centering
\caption{Modal solver vs.\ CVODE across representative cases.}
\label{tab:modal-vs-cvode}
\begin{tabular}{@{}llll@{}}
\toprule
Case & Modal & CVODE & Speedup \\
\midrule
Constant (coupled\_scalars, $N=256$) & 0.003\,s & 4.27\,s & $\mathbf{1{,}451\times}$ \\
Constant (Gertsenshtein, $N=512$) & 0.02\,s & 1.66\,s & $\mathbf{83\times}$ \\
Position-dep.\ (localized Gertsenshtein, $N=512$) & 1.4\,s & ${\sim}120$\,s & $\mathbf{{\sim}86\times}$ \\
Constraint (coupled\_proca\_3d, $16^3$) & 0.08\,s & 19.0\,s & $\mathbf{{\sim}238\times}$ \\
Long runs (Gertsenshtein, $t_\text{end}=500$) & 1.6\,s & 12.2\,s & $\mathbf{7.7\times}$ \\
\bottomrule
\end{tabular}
\end{table}

\subsection{Auto-Selection}
\label{sec:modal-auto-selection}

The modal solver is auto-selected (step~2 in \texttt{\_resolve\_scheme()}) when \emph{all} conditions are met:

\begin{enumerate}
  \item \textbf{Flat metric} --- \texttt{spec.canonical.volume\_element is None} (checked first; curved metrics like spherical/cylindrical have non-None volume elements)
  \item \textbf{No constraint equations} --- no \texttt{time\_order=0} equations
  \item \textbf{All BCs periodic} --- \texttt{grid.periodic} all \texttt{True}
  \item \textbf{All operators supported} --- all RHS operators in \texttt{\_EXACT\_MULTIPLIERS}
  \item \textbf{No time-dependent coefficients} --- \texttt{term.time\_dependent} all \texttt{False}
\end{enumerate}

Position-dependent coefficients are OK (handled via convolution). The full auto-selection hierarchy:

\begin{lstlisting}[style=bash]
1. Modal eligible (flat+periodic+time-independent+supported operators)
   -> modal (handles constraints via Schur complement if present)
2. Constraints (time_order=0) not modal-eligible -> IDA
3. First-order (time_order=1)         -> IDA
4. Dissipation (first_derivative_t)   -> IDA
5. Missing canonical (wave eq)        -> warning + CVODE
6. Pure wave, Hamiltonian             -> CVODE
\end{lstlisting}

Users can explicitly request \texttt{--scheme modal} (validated for eligibility) or override with \texttt{--scheme cvode}.

\subsubsection{Position-Dependent + Periodic: Root Cause and Fix (GH \#367)}
\label{sec:modal-pos-dep-auto-route}

\textbf{Resolved in v0.41.6 (2026-05-19).} Theories with
position-dependent coefficients and a non-trivial
\texttt{kinetic\_coefficient\_symbolic} (Gertsenshtein graviton modes
carry $-1/\kappa^2$; massive-vector torsion carries $-\xi$) previously
produced catastrophically wrong answers under the modal solver:
\texttt{expm\_multiply} reported amplitudes $\sim 10^9 \times$ too large
on the E.0 dual-Gaussian reproducer at $t_{\mathrm{end}} = 20$.

\paragraph{Root cause (verified 2026-05-19).} The two convolution-free
matrix builders in \texttt{modal.py} had drifted out of sync on the
$M^{-1}$ (inverse kinetic coefficient) folding:
\begin{itemize}
  \item \texttt{\_build\_per\_mode\_matrices} (\texttt{modal.py}:1318+)
    folds \texttt{build\_inverse\_kinetic\_diag} into the velocity-row
    coefficient at line 1411, so the emitted matrix represents
    $dv/dt = M^{-1} \cdot K(q)$.
  \item \texttt{\_build\_convolution\_matrix} (\texttt{modal.py}:1454+)
    did \emph{not} apply this scaling. For an equation with
    $\text{kinetic\_coefficient} = -1/\kappa^2$, the discrete
    Laplacian's diagonal at the Nyquist mode was
    $(-1/\kappa^2) \cdot (-k_{\max}^2) = +k_{\max}^2 / \kappa^2$ — the
    \emph{wrong sign}. The resulting 2$\times$2 state block at every
    mode $m$ had real eigenvalues $\pm k_m / \kappa$, not the
    physically-correct imaginary $\pm i k_m / \kappa$.
\end{itemize}
The visible symptom looked like a non-normal Fourier-convolution
pseudospectral artifact (max $\Re(\lambda) = k_{\max}$ tracking the
Nyquist wavenumber exactly), so an early diagnosis attributed it to
``the column placement of $i k_{m'}$ in $A[m, m'] = \widehat{c}(m-m')
\cdot i k_{m'}$''. The empirical evidence ruled that out: an mpmath
oracle at \texttt{dps=50} gave the same wildly-wrong answer as float64,
proving the matrix itself was wrong (not the precision); zeroing every
off-diagonal block left $\max\Re(\lambda)$ unchanged, locating the
bad eigenvalue on the block-diagonal Laplacian contribution rather than
the cross-coupling; and the artifact persisted at $B_{\mathrm{peak}}=0$
(no position-dependent variation at all), which definitively pointed
to a constant-coefficient sign error.

\paragraph{Fix.} \texttt{\_build\_convolution\_matrix} now calls
\texttt{build\_inverse\_kinetic\_diag(spec, params)} once at the top and
threads a per-equation \texttt{scale} kwarg through
\texttt{\_add\_convolution\_coupling}, so every velocity-row
contribution (constant-coefficient and position-dependent alike) is
multiplied by $M_{ii}^{-1}$. \texttt{build\_inverse\_kinetic\_diag}
returns \texttt{None} when every field has $M \approx 1$, so theories
without non-trivial kinetic coefficients are unaffected.

\paragraph{In-process timing on the E.0 reproducer (2026-05-19 fix).}
At $L = 100$, dual-Gaussian background, $t_{\mathrm{end}} = 20$:
\begin{center}
\begin{tabular}{rrrrr}
\toprule
$N$ & modal & CVODE & IDA & h$_5$ err vs CVODE \\
\midrule
64 & 29 ms & 86 ms & 245 ms & 0.47\% \\
128 & 48 ms & 364 ms & 1466 ms & 0.01\% \\
256 & 130 ms & 29 ms & 47 ms & 0.00\% \\
\bottomrule
\end{tabular}
\end{center}
The modal path is correct, $\sim$3$\times$ faster than CVODE at
$N=64$, and $\sim$8$\times$ faster than IDA. At $N=256$ the dense
convolution matrix grows costly enough that CVODE/IDA win; the modal
crossover is around $N \approx 200$.

\paragraph{Auto-route is now a no-op.} The
\texttt{\_override\_pos\_dep\_periodic\_scheme} helper in
\texttt{tidal/cli/\_simulate.py} is retained as a stub (returns
\texttt{(scheme, None)}) for backward-compatible signatures, but no
longer routes anywhere — modal handles pos-dep + periodic theories
correctly in all the cases the override previously diverted. The
post-evolution amplitude-growth check in
\texttt{\_evolve\_full\_matrix} (threshold $10^6\times$ initial
amplitude) stays as a defensive safety net for any future divergence.

\paragraph{What was checked and rejected during the false diagnosis.}
Recorded so future sessions don't re-investigate:
\begin{itemize}
  \item Arnoldi-Krylov / Sidje-Expokit / KIOPS as replacements for
    \texttt{expm\_multiply}: a hand-coded Sidje 1998 implementation
    ($\sim$470 LOC, adaptive sub-stepping, augmented Hessenberg error
    estimator, DGKS reorthogonalisation) gave the same wrong answer as
    dense Padé to 1$\times 10^{-14}$. The matrix itself was wrong,
    so any algorithm computing $\exp(A t) y_0$ accurately on the
    \emph{wrong} matrix returns the same wrong answer.
  \item Hou-Li exponential filter on $\widehat{c}(q)$: rejected — the
    artifact is on the block-diagonal Laplacian, not the
    cross-coupling convolution kernel.
  \item Skew-symmetric projection $A_{\mathrm{skew}} = \frac{1}{2}(A -
    A^\dagger)$: stabilised the spectrum (purely imaginary) but
    converged to a discretization of the wrong continuous operator
    (72\% error at every $N$). It dropped a Hermitian correction term
    that the physics actually needs.
  \item Row-k recombination of identity-$B'$ + gradient-$B$ pairs in
    the cross-coupling block: no change to the spectrum; confirmed
    via per-block eigenvector analysis that the dominant eigenvector
    lives on the block-diagonal (slot $\to$ velocity\_slot) Laplacian
    contribution, not the cross-coupling.
\end{itemize}
The deeper lesson is that the "non-normal pseudospectral phenomenon"
narrative was plausible but wrong: \emph{the matrix had a real
eigenvalue at every $k$, not just at the high-$k$ Nyquist edge}, and
that was the diagnostic signature that finally pointed to the missing
$M^{-1}$ scaling. The Nyquist concentration of the dominant
eigenvector was an artifact of where $k^2$ is largest, not a property
of the underlying bug.

\subsection{Usage}
\label{sec:modal-usage}

\begin{lstlisting}[style=bash]
# Auto-selected for eligible systems:
tidal simulate examples/data/coupled_scalars.json \
  --periodic --t-end 10

# Explicit selection:
tidal simulate examples/data/coupled_scalars.json \
  --scheme modal --periodic --t-end 10

# In sweeps:
tidal sweep examples/data/coupled_scalars.json \
  --sweep "gCpl=0.1:1.0:20" \
  --measure conversion --scheme modal
\end{lstlisting}

\subsection{First-Order Systems (Diffusion)}
\label{sec:modal-first-order-systems}

The modal solver handles first-order equations naturally:
\begin{equation}
  \partial_t u = D \cdot \partial_x^2 u \quad\to\quad \text{per-mode: } \frac{\mathrm{d}\hat{u}_k}{\mathrm{d}t} = -D k^2 \hat{u}_k
\end{equation}

Solution: $\hat{u}_k(t) = \hat{u}_k(0) \cdot \exp(-D k^2 t)$ --- exponential decay per mode, diagonal matrix.

\subsection{Error Model}
\label{sec:modal-error-model}

The dominant error source is the Pad\'{e} scaling-and-squaring backward error (${\sim}10^{-14}$ per application; Higham 2009 \S4) plus $O((k-1) \cdot u \cdot \|\exp(M \cdot \mathrm{d}t)\|^2)$ accumulated matmul roundoff for snapshot $k$ in the step-by-step path (see \cref{sec:modal-snapshot-precision}). For inference runs using $N=2$ snapshots (initial + final), the final snapshot is produced by a single matrix-exponential application with zero accumulation --- machine-precision at $t_\mathrm{end}$. For parameter sweeps, numerical uncertainty is effectively zero in all configurations --- the dominant error becomes the physics model itself.

Comparison with other solvers on single-field KG ($k=1$, $m^2=1$, $N=64$, $t=1$):

\begin{table}[htbp]
\centering
\caption{Error comparison: modal solver vs.\ CVODE for single-field Klein--Gordon.}
\label{tab:modal-error}
\begin{tabular}{@{}lll@{}}
\toprule
Solver & Error vs analytical & Note \\
\midrule
Modal & $5\times 10^{-16}$ & Machine precision \\
CVODE (\texttt{rtol}=$10^{-10}$) & $2.8\times 10^{-4}$ & Truncation error accumulates \\
\bottomrule
\end{tabular}
\end{table}

\subsubsection{Performance Benchmarks}
\label{sec:modal-perf-benchmarks}

Coupled scalars (2 fields, 31 snapshots, $t_\text{end}=3$, periodic BCs):

\begin{table}[htbp]
\centering
\caption{Modal vs.\ CVODE performance for coupled scalars.}
\label{tab:modal-coupled-perf}
\begin{tabular}{@{}lcccc@{}}
\toprule
Grid & Modal & CVODE (\texttt{rtol}=$10^{-10}$) & Speedup & Max diff \\
\midrule
$N=64$ & 0.003\,s & 0.16\,s & $\mathbf{57\times}$ & $4.0\times 10^{-4}$ \\
$N=128$ & 0.003\,s & 0.40\,s & $\mathbf{157\times}$ & $5.2\times 10^{-4}$ \\
$N=256$ & 0.003\,s & 4.27\,s & $\mathbf{1{,}451\times}$ & $6.3\times 10^{-4}$ \\
\bottomrule
\end{tabular}
\end{table}

Key observations:
\begin{itemize}
  \item \textbf{Modal time is $O(1)$ in grid size} --- Pad\'{e} precomputation cost is per-mode ($4\times 4$ matrices), independent of $N$. The IFFT reconstruction is $O(N \log N)$ but negligible.
  \item \textbf{CVODE scales as $O(N^{2+})$} --- spatial operator evaluation is $O(N)$ per timestep, and more grid points require more timesteps to resolve the same physics.
  \item \textbf{Max diff is CVODE error, not modal error} --- modal solutions are exact to machine precision (${\sim}10^{-16}$); the $10^{-4}$ difference is entirely CVODE truncation error.
  \item \textbf{Speedup grows with $N$} --- at $N=256$ the modal solver is ${>}1000\times$ faster, making it transformative for parameter sweeps where hundreds of simulations are needed.
  \item \textbf{Modal time is $O(1)$ in simulation time} --- $\exp(M \cdot \mathrm{d}t)$ is precomputed once per mode; each additional snapshot is one matvec at $O(bs^2)$. CVODE must take more timesteps for longer runs. See scaling table below.
\end{itemize}

\subsubsection{Gertsenshtein Effect (6 Fields, Gradient Coupling)}
\label{sec:modal-gertsenshtein-bench}

$\Bzero$ sweep ($N=256$, $L=100$, $t_\text{end}=50$, $k=2.01$, $\kappa=1$, 10 points), $P_\text{final}$ compared to analytical $\Pconv = \sin^2(\kappa \Bzero D/2) \times k^2/(k^2 + \kappa^2 \Bzero^2)$:

\begin{table}[htbp]
\centering
\caption{Gertsenshtein $\Bzero$ sweep: modal vs.\ CVODE vs.\ analytical.}
\label{tab:modal-gertsenshtein}
\begin{tabular}{@{}lcccc@{}}
\toprule
$\Bzero$ & $P$(modal) & $P$(CVODE) & $P$(analytical) & Error \\
\midrule
0.010 & 0.0612 & 0.0612 & 0.0612 & $2\times 10^{-6}$ \\
0.031 & 0.4924 & 0.4924 & 0.4923 & $1\times 10^{-4}$ \\
0.052 & 0.9313 & 0.9312 & 0.9307 & $6\times 10^{-4}$ \\
0.073 & 0.9325 & 0.9325 & 0.9314 & $1\times 10^{-3}$ \\
0.094 & 0.4949 & 0.4948 & 0.4940 & $9\times 10^{-4}$ \\
0.116 & 0.0624 & 0.0624 & 0.0623 & $5\times 10^{-5}$ \\
0.137 & 0.0740 & 0.0740 & 0.0734 & $6\times 10^{-4}$ \\
0.158 & 0.5180 & 0.5180 & 0.5143 & $4\times 10^{-3}$ \\
0.179 & 0.9435 & 0.9435 & 0.9360 & $7\times 10^{-3}$ \\
0.200 & 0.9181 & 0.9181 & 0.9105 & $8\times 10^{-3}$ \\
\bottomrule
\end{tabular}
\end{table}

\begin{itemize}
  \item \textbf{Both solvers agree to ${\sim}10^{-5}$} with each other --- identical physics, different numerics
  \item \textbf{RMS error vs analytical: 0.0036} (0.36\%) --- dominated by the effective-mass correction approximation, not solver error
  \item \textbf{Modal: 0.49\,s/point, CVODE: 1.66\,s/point $\to$ $3.4\times$ speedup} at $N=256$
  \item Speedup is lower than coupled\_scalars because the 6-field Gertsenshtein system has $12\times 12$ per-mode matrices (vs $4\times 4$ for 2 fields) and $6\times$ more IFFT reconstructions per snapshot
\end{itemize}

\subsubsection{Scaling with Simulation Time ($t_\text{end}$)}
\label{sec:modal-scaling-tend}

Gertsenshtein (6 fields, $N=256$, 101 snapshots, $\Bzero=0.1$):

\begin{table}[htbp]
\centering
\caption{Modal vs.\ CVODE scaling with simulation time.}
\label{tab:modal-scaling-tend}
\begin{tabular}{@{}lccc@{}}
\toprule
$t_\text{end}$ & Modal & CVODE & Speedup \\
\midrule
10  & 1.4\,s & 1.7\,s  & $\mathbf{1.2\times}$ \\
50  & 1.5\,s & 3.1\,s  & $\mathbf{2.0\times}$ \\
100 & 1.6\,s & 4.3\,s  & $\mathbf{2.8\times}$ \\
200 & 2.1\,s & 7.2\,s  & $\mathbf{3.4\times}$ \\
500 & 1.6\,s & 12.2\,s & $\mathbf{7.7\times}$ \\
\bottomrule
\end{tabular}
\end{table}

\begin{itemize}
  \item \textbf{Modal cost is constant in $t_\text{end}$} (${\sim}1.5$\,s) --- $\exp(M \cdot \mathrm{d}t)$ is precomputed once per mode, and each snapshot step is one matvec at $O(bs^2)$, $O(1)$ in $\Delta t$. The ${\sim}1.5$\,s is Python startup + FFT + Pad\'{e} precomputation overhead.
  \item \textbf{CVODE cost scales linearly with $t_\text{end}$} --- each additional unit of simulation time requires proportionally more adaptive timesteps.
  \item \textbf{Speedup grows linearly with $t_\text{end}$} --- for a 40-point parameter sweep at $t_\text{end}=500$, modal saves ${\sim}7$ minutes vs CVODE. For $t_\text{end}=5000$ (astrophysical timescales), the advantage would be ${\sim}50\times$.
\end{itemize}

\subsection{Implementation Details}
\label{sec:modal-implementation}

\subsubsection{Files}
\label{sec:modal-files}

\begin{table}[htbp]
\centering
\caption{Modal solver source files.}
\label{tab:modal-files}
\begin{tabular}{@{}ll@{}}
\toprule
File & Purpose \\
\midrule
\texttt{tidal/solver/modal.py} & Core solver (${\sim}1{,}400$ lines) \\
\texttt{tidal/cli/\_simulate.py} & Auto-selection + dispatch \\
\texttt{tidal/cli/\_\_init\_\_.py} & \texttt{"modal"} in \texttt{--scheme} choices \\
\texttt{tests/test\_solver\_modal.py} & 35 tests (12 eligibility + 14 correctness + 4 stability + 5 constraint) \\
\bottomrule
\end{tabular}
\end{table}

\subsubsection{Key Functions}
\label{sec:modal-key-functions}

\begin{table}[htbp]
\centering
\caption{Modal solver key functions.}
\label{tab:modal-functions}
\begin{tabular}{@{}ll@{}}
\toprule
Function & Purpose \\
\midrule
\texttt{can\_use\_modal(spec, grid, bc)} & Eligibility check (shared with auto-selection) \\
\texttt{solve\_modal(spec, grid, y0, \ldots)} & Main entry point --- routes to per-mode, full-matrix, or unified-reduction path \\
\texttt{\_fft\_slots / \_ifft\_slots} & State $\leftrightarrow$ Fourier transforms (analytical shape computation) \\
\texttt{\_build\_per\_mode\_matrices} & Per-mode evolution matrices (constant coefficients, no constraints) \\
\texttt{\_build\_convolution\_matrix} & Convolution matrix with vectorized coupling (position-dependent) \\
\texttt{\_build\_evolution\_matrices} & Unified builder: algebraic constraints (Schur) + rank-deficient mass-matrix (eigen-Schur) + near-singular vel\_coupling. Returns $(A_{\text{rhs}}, B_{\text{lhs}})$ for generalized eigenvalue solve \\
\texttt{\_evolve\_per\_mode} & Block-aware eigendecomposition solver; optional Fourier snapshot return \\
\texttt{\_evolve\_full\_matrix} & Sparse-aware \texttt{expm\_multiply} solver (auto-thresholds to CSC) \\
\texttt{\_find\_independent\_blocks} & Union-find block detection for eigendecomposition \\
\bottomrule
\end{tabular}
\end{table}

\subsubsection{Reused Infrastructure}
\label{sec:modal-reused}

\begin{table}[htbp]
\centering
\caption{Infrastructure reused by the modal solver.}
\label{tab:modal-reused}
\begin{tabular}{@{}lll@{}}
\toprule
Component & Location & How Used \\
\midrule
\texttt{\_get\_wavenumbers()} & \texttt{operators.py} & $k$-grid construction \\
\texttt{is\_periodic\_bc()} & \texttt{operators.py} & Eligibility gate \\
\texttt{StateLayout} & \texttt{state.py} & Field$\leftrightarrow$state mapping \\
\texttt{CoefficientEvaluator} & \texttt{coefficients.py} & Coefficient resolution \\
\texttt{SimulationProgress} & \texttt{progress.py} & tqdm progress bar \\
\bottomrule
\end{tabular}
\end{table}

\textbf{Not reused:} \texttt{constraint\_solve.py:\_MULTIPLIERS} --- uses modified wavenumbers matching FD stencils. Modal solver defines its own exact multiplier registry.

\subsection{Kinetic Coefficient Handling --- Two Dispatches}
\label{sec:modal-kinetic-dispatches}

The mass matrix $M$ from \texttt{kinetic\_coefficient\_symbolic} (``$M \ddot{x} = K x$'' rather than ``$\ddot{x} = K x$'') reaches the solver through \emph{two distinct code paths} depending on whether the spec has algebraic constraints or higher-order time operators. Both paths must apply $M$ identically, or cross-path regressions appear silently (the original \#301 Bug B symptom). This section documents the contract so future changes update both paths.

\textbf{Why two paths.} The fast path exists because the majority of theories (Klein--Gordon, wave equations, most multi-field systems) have neither constraints nor higher-order time operators on the RHS. For these, bypassing the generalized eigenvalue problem and producing the pre-solved first-order evolution matrix directly saves the QZ decomposition cost and avoids spurious infinite eigenvalues that appear when $B$ is exactly $I$.

\begin{table}[htbp]
\centering
\caption{Kinetic coefficient consumption across modal code paths. The dispatch condition is computed by \texttt{solve\_modal} at line 2385 as \texttt{needs\_reduction = (has\_constraints or has\_time\_ops) and not has\_pos\_dep}.}
\label{tab:modal-kinetic-dispatch}
\begin{tabular}{@{}p{0.18\textwidth}p{0.32\textwidth}p{0.42\textwidth}@{}}
\toprule
Path & Dispatch condition & How $M$ is consumed \\
\midrule
Fast (\texttt{\_build\_per\_mode\_matrices}) & No constraints, no d\_t RHS ops, no pos.\ dep.\ coefficients & \texttt{build\_inverse\_kinetic\_diag(spec, params)} computes $M^{-1}$ per field; folded into velocity-row coefficients so \texttt{A[vel\_slot, target\_slot] = $M^{-1}$ · coeff · multiplier}. Single eigendecomposition of the resulting first-order matrix. \\
Gen.\ eig (\texttt{\_build\_evolution\_matrices}) & Constraints, or d\_t RHS ops, or gauge $A_0$ demotion & \texttt{kinetic\_coefficient\_symbolic} populated on the diagonal of \texttt{M\_mat}; \texttt{scipy.linalg.eig(A, B)} with QZ handles $M$, rank-deficiency via null-space projection, hidden algebraic constraints via Schur elimination. \\
\bottomrule
\end{tabular}
\end{table}

\textbf{Shared helper.} \texttt{tidal/solver/\_kinetic.py::build\_inverse\_kinetic\_diag(spec, params)} is the single entry point consumed by:
\begin{itemize}
  \item Modal fast path (this path --- folds $M^{-1}$ into the evolution matrix)
  \item \texttt{solver/scipy\_solver.py}, \texttt{solver/cvode.py}, \texttt{solver/leapfrog.py}, \texttt{solver/ida.py} (time-domain backends --- multiply RHS by $M^{-1}$ at each step; IDA absorbs it into the DAE residual)
\end{itemize}
The modal generalized-eig path is the exception: it reads \texttt{kinetic\_coefficient\_symbolic} directly into \texttt{M\_mat}, because the generalized eigenvalue solve needs $M$ as a matrix, not as a scalar-per-field inverse. These are two representations of the same physical input.

\textbf{Fast-path trigger.} \texttt{build\_inverse\_kinetic\_diag} returns \texttt{None} when every dynamical field has $|M_{ii} - 1| < 10^{-12}$, letting every caller skip the per-step multiply. Theories without \texttt{kinetic\_coefficient\_symbolic} (the vast majority) pay zero overhead for this feature.

\textbf{Sign convention.} \texttt{kinetic\_coefficient\_symbolic} carries the raw LHS coefficient sum from the Wolfram EOM (see \cref{sec:pr-eh-cd}), which may be negative (e.g., graviton $-\kappa^{-2}$). Both paths accept negative $M$ --- negative eigenvalues of the generalized problem give ghost modes, which are physically meaningful for theories like PGT with indefinite kinetic structure. The consistency test \texttt{tests/test\_solver\_kinetic\_consistency.py} validates agreement across all five backends on a theory with $M = 1.5$.

\textbf{Regression guard.} Both paths must be updated together when the kinetic handling changes. The cross-backend consistency test will catch any path that stops consuming $M$ correctly.

\textbf{Scope.} Current implementation handles \emph{diagonal} $M$ (one scalar per field). Off-diagonal kinetic mixing (e.g., $\delta_m F \cdot F_{\text{torsion}}$ in dark-photon plasma theories) is a separate extension tracked in issue \#305; the data structures in both paths are already $(n_{\text{fields}}, n_{\text{fields}})$-shaped so the diagonal case is a 1D slice of the general case.

\subsection{Block-Aware Eigendecomposition}
\label{sec:modal-block-aware}

Multi-field systems often have block-diagonal per-mode matrices (e.g.\ Gertsenshtein: $h_5 \leftrightarrow a_1$ and $h_7 \leftrightarrow a_2$ as independent $4\times 4$ blocks). The modal solver detects these blocks via union-find on the coupling graph and eigendecomposes each independently. This:

\begin{enumerate}
  \item \textbf{Prevents degenerate-eigenvalue mixing} --- \texttt{np.linalg.eig} on a full matrix with repeated eigenvalues across independent blocks can mix eigenvectors, projecting nonzero-IC components onto zero-IC blocks
  \item \textbf{Enables zero-IC block skipping} --- blocks where all initial conditions are zero produce exact zeros without eigendecomposition
  \item \textbf{Reduces computation} --- many small eigendecompositions instead of one large one
\end{enumerate}

For systems with gradient coupling (e.g.\ Gertsenshtein $h \leftrightarrow a$), low-$k$ modes may have genuinely positive real eigenvalues (physical parametric conversion). A warning is issued when $\max(\mathrm{Re}(\lambda)) \times \Delta t > 30$ (approaching overflow).

\subsection{Constraint Elimination (Fourier Schur Complement)}
\label{sec:modal-constraint-elimination}

Systems with constraint equations (\texttt{time\_order=0}, e.g., Proca $A_0$) are handled via algebraic elimination in Fourier space. The constraint self-operator $L$ (e.g., $m^2 - \nabla^2$) has exact Fourier multipliers, so $L^{-1}$ is trivially $1/(m^2 + k^2)$ per mode.

\textbf{How it works}: For a mixed system with dynamical ($d$) and constraint ($c$) fields:
\begin{align}
  \frac{\mathrm{d}}{\mathrm{d}t}[d] &= A_{dd} \cdot d + A_{dc} \cdot c & &\text{(dynamical equations)} \label{eq:schur-dyn} \\
  0 &= S_{cd} \cdot d + S_{cc} \cdot c & &\text{(constraint defines $c$ algebraically)} \label{eq:schur-con}
\end{align}

Substituting $c = -S_{cc}^{-1} \cdot S_{cd} \cdot d$ gives the reduced system:
\begin{equation}
  \frac{\mathrm{d}}{\mathrm{d}t}[d] = \bigl(A_{dd} - A_{dc} \cdot S_{cc}^{-1} \cdot S_{cd}\bigr) \cdot d
  \label{eq:schur-reduced}
\end{equation}

The constraint velocity coupling (when dynamical equations reference $v_{A_0} = \partial_t A_0$) creates an implicit equation resolved by batched \texttt{np.linalg.solve} across all modes: $(I - A_{dc,\text{vel}} \cdot S_{cc}^{-1} \cdot S_{cd}) \cdot d' = A_\text{rhs} \cdot d$.

\textbf{Optimizations:}
\begin{itemize}
  \item \textbf{Batched $S_{cc}$ inversion}: \texttt{np.linalg.inv(S\_cc)} inverts all modes at once (no Python loop), with vectorized singularity detection for gauge-freedom modes ($k=0$)
  \item \textbf{Batched implicit solve}: \texttt{np.linalg.solve(lhs, A\_rhs)} across all modes simultaneously
  \item \textbf{Fourier snapshot reuse}: \texttt{\_evolve\_per\_mode(return\_fourier=True)} saves the already-computed Fourier-space data at each snapshot, eliminating $n_\text{snapshots} \times n_\text{dyn}$ redundant forward FFTs during constraint field reconstruction
\end{itemize}

At each output time, constraint fields are reconstructed directly from the saved Fourier data: $c(k) = -S_{cc}^{-1} \cdot S_{cd} \cdot d(k)$.

\textbf{Performance}: On coupled\_proca\_3d ($16^3$, $t=5$, 51 snapshots): modal 0.08\,s vs IDA 19.0\,s --- $\mathbf{{\sim}238\times}$ \textbf{speedup}. IDA fails entirely at $32\times 32$ (convergence issues), while modal runs at any grid size.

\textbf{Automatic constraint IC satisfaction}: Unlike non-modal solvers, which rely on a separate \texttt{ensure\_consistent\_ic()} call in \texttt{constraint\_solve.py} to pre-solve constraint fields before time integration, the modal solver's Schur complement provides automatic constraint IC satisfaction as part of the elimination step. Constraint field values are reconstructed algebraically from the dynamical amplitudes ($c = -S_{cc}^{-1} \cdot S_{cd} \cdot d$) at every output time, so the initial conditions are consistent by construction --- no separate constraint solver invocation is needed.

\textbf{References}: Hairer \& Wanner (1996), \emph{Solving ODEs II}, Ch.\ VII; Ascher \& Petzold (1998), \emph{Computer Methods for ODEs/DAEs}, \S10.2.

\subsection{Generalized Mass Matrix (Higher-Derivative Theories)}
\label{sec:modal-mass-matrix}

Higher-derivative theories such as PGT with $\tilde{R}^2$ terms produce equations with implicit acceleration coupling: \texttt{d2\_t(field\_j)} appearing on the RHS of \texttt{field\_i}'s equation, meaning $\ddot{q}_j$ couples into $\ddot{q}_i$'s dynamics. After v6 perturbative reduction (\cref{sec:perturbative-reduction}; see also the architectural Ostrogradsky discussion in \cref{sec:ostrogradsky}), the system has the generalized form:
\begin{equation}
  M(k) \cdot \ddot{x} = K(k) \cdot x + D(k) \cdot \dot{x} + J(k) \cdot \dddot{x}
  \label{eq:generalized-mass}
\end{equation}
where $M$ is the \emph{mass matrix} (from \texttt{d2\_t}, \texttt{mixed\_T2\_S*} operators), $K$ is the stiffness matrix (pure spatial operators), $D$ is the damping matrix (\texttt{first\_derivative\_t}, \texttt{mixed\_T1\_S*} operators), and $J$ is the jerk matrix (\texttt{d3\_t}, \texttt{mixed\_T3\_S*}).

\subsubsection{Operator Decomposition}
\label{sec:modal-operator-decomp}

Every operator in flat spacetime decomposes as $\text{spatial\_multiplier}(k) \times \partial_t^n$:
\begin{table}[htbp]
\centering
\caption{Operator decomposition into spatial multiplier and time derivative order.}
\label{tab:operator-decomposition}
\begin{tabular}{@{}llcl@{}}
\toprule
Operator & Spatial mult.\ $\hat{O}(k)$ & Time order $n$ & Matrix \\
\midrule
\multicolumn{4}{@{}l}{\emph{Position operators ($n=0$): contribute to $K(k)$}} \\
\texttt{identity} & $1$ & 0 & $K$ \\
\texttt{laplacian} / \texttt{laplacian\_x/y/z} & $-k_x^2$ (or $-|\mathbf{k}|^2$) & 0 & $K$ \\
\texttt{gradient\_x/y/z} & $ik_x$ & 0 & $K$ \\
\texttt{cross\_derivative\_xy/xz/yz} & $-k_x k_y$ & 0 & $K$ \\
\texttt{biharmonic} & $|\mathbf{k}|^4$ & 0 & $K$ \\
\texttt{derivative\_3\_x/y/z} & $-ik_x^3$ & 0 & $K$ \\
\midrule
\multicolumn{4}{@{}l}{\emph{Velocity operators ($n=1$): contribute to $D(k)$}} \\
\texttt{first\_derivative\_t} & $1$ & 1 & $D$ \\
\texttt{mixed\_T1\_S1x/y/z} & $ik_x$ & 1 & $D$ \\
\midrule
\multicolumn{4}{@{}l}{\emph{Acceleration operators ($n=2$): contribute to $M(k)$}} \\
\texttt{d2\_t} & $1$ & 2 & $M$ \\
\texttt{mixed\_T2\_S1x/y/z} & $ik_x$ & 2 & $M$ \\
\texttt{mixed\_T2\_S2x/y/z} & $-k_x^2$ & 2 & $M$ \\
\midrule
\multicolumn{4}{@{}l}{\emph{Jerk operators ($n=3$): contribute to $J(k)$, eliminated via EOM}} \\
\texttt{d3\_t} & $1$ & 3 & $J$ \\
\texttt{mixed\_T3\_S1x/y/z} & $ik_x$ & 3 & $J$ \\
\bottomrule
\end{tabular}
\end{table}

The registry contains 35 operators total (including axis variants $x/y/z$).  Each operator $\mathcal{O}$ acting on a field $\phi$ in Fourier space becomes:
\begin{equation}
  \widehat{\mathcal{O}[\phi]}(k) = \hat{O}(k) \cdot \partial_t^n \hat{\phi}(k)
\end{equation}
where $\hat{O}(k)$ is the spatial multiplier from the table and $n$ is the time order.  The sign convention follows from the Fourier derivative property: $\widehat{\partial_x \phi} = ik_x \hat\phi$, $\widehat{\partial_x^2 \phi} = -k_x^2 \hat\phi$, $\widehat{\partial_x^3 \phi} = -ik_x^3 \hat\phi$.  Mixed time-space operators (e.g., $\partial_t^2 \partial_x$) decompose as $ik_x \times \partial_t^2$: the spatial part enters the multiplier, the temporal part determines which matrix ($M$, $D$, $K$, or $J$) receives the contribution.

\subsubsection{Singular Mass Matrix}
\label{sec:modal-singular-mass}

For PGT with $\tilde{R}^2$, the mass matrix is \emph{singular}. The graviton-torsion system has $M = \bigl[\begin{smallmatrix} 1 & 1 \\ 1 & 1 \end{smallmatrix}\bigr]$ for the $\{t_3, t_7\}$ pair (eigenvalues $\{0, 2\}$) and similarly for $\{w_{h_3}, w_{h_5}\}$. The zero eigenvalue creates a hidden algebraic constraint: $t_3 = t_7$.

The algorithm eigendecomposes $M$ per mode:
\begin{enumerate}
  \item Compute $M = Q \Lambda Q^T$ (real symmetric eigendecomposition)
  \item Zero eigenvalues $\to$ constraint directions: $0 = \tilde{K}_c \cdot z + \tilde{D}_c \cdot \dot{z}$ where $z = Q^T x$
  \item Nonzero eigenvalues $\to$ dynamical directions: $\Lambda_d \cdot \ddot{z}_d = \tilde{K}_d \cdot z + \tilde{D}_d \cdot \dot{z}$
  \item Schur-eliminate constrained modes (same algorithm as \cref{sec:modal-constraint-elimination})
  \item Invert $\Lambda_d$ (diagonal, invertible) to get $\ddot{z}_d = \Lambda_d^{-1} \tilde{K}_d \cdot z + \Lambda_d^{-1} \tilde{D}_d \cdot \dot{z}$
\end{enumerate}

\subsubsection{Jerk Substitution}
\label{sec:modal-jerk}

After mass-matrix elimination, jerk coupling $J(k) \cdot \dddot{x}$ may remain (from \texttt{d3\_t} and \texttt{mixed\_T3\_S*} operators).  These arise in $\tilde{R}^2$ theories where the Ostrogradsky-reduced equations still contain third time derivatives of certain fields.  The jerk is eliminated algebraically by differentiating the (now-invertible) dynamical equations.

\paragraph{Algorithm.}
Define the mass-inverted dynamics matrices:
\begin{equation}
  E = M_{\text{dyn}}^{-1} \cdot K, \qquad F = M_{\text{dyn}}^{-1} \cdot D
\end{equation}
so that $\ddot{x} = E \cdot x + F \cdot \dot{x}$.  Then $\dddot{x} = \dot{(\ddot{x})} = E \cdot \dot{x} + F \cdot \ddot{x} = E \cdot \dot{x} + F \cdot (E \cdot x + F \cdot \dot{x})$, giving:
\begin{equation}
  \dddot{x} = F \cdot E \cdot x + (E + F^2) \cdot \dot{x}
\end{equation}
Substituting into the original system $M\ddot{x} = K x + D \dot{x} + J \dddot{x}$:
\begin{align}
  K_{\text{eff}} &= K + J \cdot F \cdot E  &\text{(position correction)} \\
  D_{\text{eff}} &= D + J \cdot (E + F^2)  &\text{(velocity correction)}
\end{align}
The implementation checks $\max|J| > 10^{-15}$ before applying the substitution; if no jerk terms are present, no correction is needed.

\paragraph{Constraint acceleration jerk.}
Constraint equations may also contain jerk operators on dynamical fields (e.g., \texttt{mixed\_T3\_S1x(t\_3)}).  These are handled separately: after mass inversion, the constraint's jerk term $\hat{O}(k) \cdot \dddot{t}_3$ is replaced using the same substitution formula, converting it to effective position and velocity terms in the constraint matrix $S_{cd}$ before the standard Schur complement.

\subsubsection{Constraint Acceleration Terms}
\label{sec:modal-constraint-accel}

Constraint equations may contain acceleration operators on dynamical fields (e.g., \texttt{mixed\_T2\_S1x(t\_3)} = $ik_x \cdot \ddot{t}_3$ in the $h_1$ constraint). After mass-matrix inversion, $\ddot{t}_3 = \sum_j E_{t_3,j} x_j + F_{t_3,j} \dot{x}_j$. Substituting into the constraint converts acceleration terms to position and velocity terms, after which the standard Schur complement applies.

\subsubsection{Velocity Coupling Singularity (QZ Decomposition)}
\label{sec:modal-qz-velocity}

When constraint fields have \emph{circular velocity dependencies} --- e.g., the $h_1$ constraint references $v_{h_0}$ while the $h_0$ constraint references $v_{h_1}$ --- the implicit velocity coupling equation
\begin{equation}
  (I - V_\text{coupling}) \cdot d' = A_\text{rhs} \cdot d
  \label{eq:vel-coupling-implicit}
\end{equation}
becomes singular ($\det(I - V_\text{coupling}) = 0$). This arises in PGT from Poincar\'e gauge symmetry, which creates hidden constraints between constraint field velocities. Na\"ive regularization (adding $\epsilon I$) produces spurious eigenvalues with $|\mathrm{Im}(\lambda)| \sim 10^8$, corrupting the evolution.

\textbf{Solution}: Reformulate as a \emph{generalized eigenvalue problem} using QZ decomposition:
\begin{equation}
  B \cdot v = \lambda \cdot A_\text{rhs} \cdot v, \qquad B = I - V_\text{coupling}
  \label{eq:qz-generalized}
\end{equation}
solved via \texttt{scipy.linalg.eig(A\_rhs, B)}. The QZ algorithm (Moler \& Stewart 1973) computes eigenvalues as ratios $\lambda_i = \alpha_i / \beta_i$:
\begin{itemize}
  \item $\beta_i \neq 0$: \textbf{Physical eigenvalue} $\lambda_i = \alpha_i/\beta_i$ --- finite, $O(1)$, governs wave dynamics
  \item $\beta_i = 0$: \textbf{Gauge mode} ($\lambda_i = \infty$) --- constrained direction from gauge symmetry. Set to zero in the evolution ($e^{0 \cdot t} = 1$), freezing gauge DOF at their initial values
\end{itemize}

\textbf{Comparison with regularization}:
\begin{center}
\begin{tabular}{@{}lcc@{}}
\toprule
Metric & Regularization ($+\epsilon I$) & QZ decomposition \\
\midrule
$\max|A|$ & $\sim 10^{19}$ & $O(1)$ \\
$\max|\mathrm{Im}(\lambda)|$ & $\sim 3.4 \times 10^8$ & $O(k^2)$ (physical) \\
Energy conservation ($t=100$) & Overflow & $\max|dE/E| = 0$ \\
\bottomrule
\end{tabular}
\end{center}

For the graviton-torsion system, circular dependencies $h_0 \leftrightarrow h_1$ and $t_0 \leftrightarrow t_8$ create gauge-constrained velocity directions at 32/33 Fourier modes. The QZ decomposition correctly identifies and filters these, producing clean $O(1)$ eigenvalues for the physical dynamics.

\textbf{References}: Golub \& Van Loan (2013), \emph{Matrix Computations}, 4th ed., \S7.7.6; Moler, C.\ \& Stewart, G.W.\ (1973), ``An Algorithm for Generalized Matrix Eigenvalue Problems'', \emph{SIAM J.\ Numer.\ Anal.}\ 10(2).

\subsubsection{Complete Elimination Hierarchy}
\label{sec:modal-elimination-hierarchy}

The generalized modal solver performs a three-level elimination, each following the same mathematical pattern (detect singular directions via eigendecomposition, project out, Schur-eliminate):

\begin{enumerate}
  \item \textbf{Level 1 --- Mass-matrix Schur}: Eigendecompose $M$ per mode. Zero eigenvalues reveal hidden constraints (e.g., $t_3 = t_7$ from Poincar\'e gauge symmetry). Schur-eliminate constrained directions; invert $M$ in dynamical subspace.
  \item \textbf{Jerk substitution}: Replace $\dddot{x}$ terms using the (now-invertible) dynamical equations: $\dddot{x}_j = E_j \cdot \dot{x} + F_j \cdot (E \cdot x + F \cdot \dot{x})$.
  \item \textbf{Constraint acceleration substitution}: Replace $\ddot{x}$ operators in constraint equations using the mass-inverted dynamics ($\ddot{x} = E \cdot x + F \cdot \dot{x}$).
  \item \textbf{Level 2 --- Constraint field Schur}: Eliminate \texttt{time\_order=0} fields via Fourier Schur complement (existing code, \cref{sec:modal-constraint-elimination}).
  \item \textbf{Level 3 --- Velocity coupling (QZ)}: Generalized eigenvalue via QZ decomposition (\cref{sec:modal-qz-velocity}). Infinite eigenvalues (gauge DOF) frozen; finite eigenvalues evolved.
\end{enumerate}

For the 4D PGT graviton-torsion system (v0.18.0, full quadratic Lagrangian), this reduces 37 fields (34 original + 3 Ostrogradsky auxiliaries) through three elimination stages:

\begin{center}
\begin{tabular}{@{}lrl@{}}
\toprule
Stage & State & Eliminated \\
\midrule
After Ostrogradsky & 37 fields, 24 dyn slots & 25 constraint fields identified \\
Level 1 (mass-matrix) & $\leq 24$ dyn slots & Hidden gauge constraints from $\det M = 0$ \\
Jerk substitution & Same dim & $\dddot{x}$ operators $\to$ $K_\text{eff}, D_\text{eff}$ \\
Constraint accel.\ sub. & Same dim & $\ddot{x}$ in constraints $\to$ $K, D$ terms \\
Level 2 (constraint Schur) & $\leq 24$ dyn slots & 25 constraint fields eliminated \\
Level 3 (QZ velocity) & Physical DOF & Gauge directions (infinite $\lambda$) frozen \\
Eigendecomposition & $\exp(\lambda \cdot t)$ & $O(1)$ time evolution \\
\bottomrule
\end{tabular}
\end{center}

Energy conservation after all three levels: $\max|dE/E| = 0$ (machine precision) at $t_\text{end} = 1.0$ for $\alpha_1 = 0.5$, $\alpha_2 = 0.3$, $\alpha_3 = 0.2$, $b_5 = 0.1$.

\textbf{Implementation}: \texttt{\_build\_evolution\_matrices()} is the single unified reduction builder --- it handles algebraic constraint fields (\texttt{time\_order=0}) via Fourier Schur complement, rank-deficient mass matrices $M$ (e.g.\ trace projections of higher-rank tensor perturbations) via per-mode mass eigendecomposition with null-space Schur elimination, and near-singular velocity coupling $B = I - V_\text{coupling}$ via the generalized eigenvalue problem. It returns $(A_\text{rhs}, B_\text{lhs})$; when $B_\text{lhs}$ is non-\texttt{None}, \texttt{\_evolve\_per\_mode()} uses \texttt{scipy.linalg.eig(A, B)} (QZ decomposition), which is numerically stable for rank-deficient and ill-conditioned $B$; otherwise it uses \texttt{np.linalg.eig(A)}. The gauge filter $|\lambda| > 10^{12}$ catches infinite eigenvalues from null directions and freezes them at IC. Prior to the \#256 unification two competing builders coexisted (constraint-eliminated and generalized); the constraint-eliminated builder silently mis-classified \texttt{d2\_t} cross-couplings as stiffness operators and was deleted --- see \cref{sec:troubleshooting-rank-deficient-kinetic} for the failure mode and its resolution.

\textbf{References}: Golub \& Van Loan (2013), \emph{Matrix Computations}, 4th ed., \S7.7; Hairer \& Lubich (2003), ZAMM 83(1); Ostrogradsky M.V.\ (1850), Mem.\ Acad.\ St.\ Petersbourg VI 4, 385.

\subsection{Constraints and Limitations}
\label{sec:modal-constraints-limitations}

\begin{enumerate}
  \item \textbf{Constraint equations} (\texttt{time\_order=0}): Handled via Fourier Schur complement (see \cref{sec:modal-constraint-elimination}) when all constraint operators have exact Fourier multipliers. Otherwise falls back to IDA.

  \item \textbf{Curved metrics} (spherical, cylindrical, polar): Not applicable --- non-periodic domains and position-dependent operators with slowly-decaying Fourier transforms create dense all-to-all coupling.

  \item \textbf{Time-dependent coefficients}: Would require ODE integration in modal space (\texttt{scipy.integrate.solve\_ivp}). Currently rejected by eligibility check (no known TIDAL examples need this).

  \item \textbf{Spectral operators and energy measurement}: The modal solver works in $k$-space, but the energy measurement evaluates the Hamiltonian in physical space. The \texttt{--spectral} flag is preserved (not disabled) for modal so that the energy measurement uses FFT-based gradient operators matching the modal solver's conserved Hamiltonian. Without this, FD gradients produce conservation errors that increase with $N$ (the FD Hamiltonian differs from the Fourier Hamiltonian). For gradient$\times$gradient terms in the Hamiltonian, the measurement uses direct gradient product $\langle \partial f, \partial g \rangle$ rather than IBP $\langle -f, \partial^2 g \rangle$ when spectral is active, because rfft Nyquist-mode handling creates $O(1/N)$ discrepancy between the two.

  \item \textbf{Per-mode eigendecomposition conditioning} (\textbf{resolved} in v0.31+; see \cref{sec:modal-robust-matrix-exp}): For coupled multi-field systems where fields have similar dispersion (e.g., massless graviton-photon, or PGT \texttt{TorsionCDT} with redundant components), the per-mode block has near-degenerate eigenvalue pairs, causing $\mathrm{cond}(V)$ up to $10^{17}$ (and up to $10^{291}$ in pathological cases). This previously caused energy-conservation degradation and path-dependent simulation results. The fix is described in the next subsection: replace the $V \cdot \mathrm{diag}(e^{\lambda t}) \cdot V^{-1}$ evolution with direct matrix-exponential evolution via Pad\'{e} scaling-and-squaring, which has backward error bounded by the function condition number rather than the eigenvector condition number.
\end{enumerate}

\subsection{Robust Matrix-Exponential Evolution}
\label{sec:modal-robust-matrix-exp}

\subsubsection{Problem: ill-conditioned eigenvector matrix \texorpdfstring{$V$}{V}}

The per-mode evolution from \cref{sec:modal-block-aware} originally computed
\begin{equation}
y(t) = V \cdot \mathrm{diag}\bigl(e^{\lambda_k t}\bigr) \cdot V^{-1} \cdot y(0)
\label{eq:modal-eigen-evolution}
\end{equation}
where $V$ is the per-mode eigenvector matrix and $\lambda_k$ are the eigenvalues of the per-mode block $A$ (or generalized eigenvalues of the pencil $(A,B)$ when \texttt{B\_modes} is non-\texttt{None}). For coupled multi-field PGT theories with redundant tensor components --- in particular, the TorsionCDT formulation of the dark-photon Lagrangian (\#318, \#320) --- the redundant degrees of freedom drive $\mathrm{cond}(V)$ up to $10^{15}\text{--}10^{17}$ at every parameter point. Per Higham 2008 Functions of Matrices, \S2.3:
\begin{quote}
``If $\kappa(V) \cdot u > 10^{-3}$, abandon diagonalization.''
\end{quote}
With IEEE-754 unit roundoff $u \approx 10^{-16}$, the abandonment threshold is $\kappa(V) > 10^{13}$. Our observed $10^{15}\text{--}10^{17}$ is two to four orders of magnitude past every published threshold. Trefethen \& Bau (1997) Lecture~33 makes the consequence concrete: at $\mathrm{cond}(V) > 10^{14}$, $V^{-1}$ has \emph{zero useful digits} in IEEE double precision. The empirical signal: identical CDT and FV (fundamental-vector, 10-field) formulations of the same physics gave inconsistent $P_\text{max}$ at the same parameter point (\(P_\text{max} = 0.571\) vs.\ \(0.063\) at $\xi = 4.5$, with FV the consistent reference per the cross-formulation cross-check).

\subsubsection{Why pseudo-inverse is the wrong fix}
\label{sec:modal-pinv-wrong}

A natural patch is to replace \texttt{np.linalg.inv} with \texttt{np.linalg.pinv} (Moore--Penrose pseudo-inverse via SVD). This is mathematically incorrect for our problem because:
\begin{enumerate}
  \item \texttt{pinv} truncates singular values below $\mathrm{rcond} \cdot \sigma_\text{max}$ ($\mathrm{rcond} \approx 10^{-13}$ in numpy 1.26+). The discarded right-singular vectors lie in directions of physical state space; there is no guarantee these directions carry no signal at our parameter points.
  \item The eigendecomposition $A = V \Lambda V^{-1}$ is \emph{not} an orthogonal decomposition (since $V$ is not unitary for non-normal $A$). SVD-based pseudo-inverse solves a different problem (minimum-norm least-squares of a singular linear system) and computes $f(A)$ only by accident when the discarded subspace happens to carry no signal.
  \item Higham 2008 \S2.3 lists pinv-of-$V$ as a workaround, not a method, and warns that using it in $f(A) = V f(\Lambda) V^{+}$ produces ``results that may bear no relation to $f(A)$'' when $V$ is near-singular.
\end{enumerate}

\subsubsection{Principled replacement: matrix exponential without diagonalization}
\label{sec:modal-pade-evolution}

The fix is to evaluate $\exp(M \cdot t) \cdot y_0$ directly where $M \equiv B^{-1} A$ (one stable solve, since $\mathrm{cond}(B) \approx 1$ for our kinetic matrix), \emph{without} forming an eigenvector basis. Two production-grade scipy routines exist:
\begin{description}
  \item[Method A: \texttt{scipy.linalg.expm}] (Higham 2009 \emph{SIAM Review} 51(4)): Pad\'{e} scaling-and-squaring. Computes the full $\exp(M)$ matrix via diagonal Pad\'{e} approximant $r_m(M / 2^s)$ followed by $s$ squaring steps. Backward error bounded by the function condition number $\kappa_f(M, t)$, not by the eigenvector condition number $\kappa(V)$.
  \item[Method B: \texttt{scipy.sparse.linalg.expm\_multiply}] (Al-Mohy \& Higham 2011): action of the matrix exponential, computing $\exp(t M) \cdot b$ directly via scaling and truncated Taylor series on the action vector. Already used in the \emph{position-dependent} path \texttt{\_evolve\_full\_matrix} (\cref{sec:modal-perf-optimizations}, Round~1) for non-normal convolution matrices.
\end{description}

Both methods are mathematically robust for arbitrary $\mathrm{cond}(V)$. Their wall-time characteristics differ substantially for the per-mode block sizes $bs \le 15$ relevant here, and the empirically-correct choice (Method A with precomputed $\exp(M\,\mathrm{d}t)$ + matvec in the snapshot loop) is non-obvious; \cref{tab:modal-expm-benchmark} below documents the benchmark.

\subsubsection{Wall-time benchmark across campaign workloads}
\label{sec:modal-expm-benchmark}

Four paths were benchmarked on a representative four-row workload table covering Stage~A (\texttt{dark\_photon\_plasma} CDT, \texttt{torsion\_dark\_photon\_fv}) and Stages~B-E (\texttt{graviton\_torsion}, \texttt{euler\_heisenberg}). Synthetic per-mode matrices match each theory's block dimension $bs$, block count, $n_\text{modes}$, $n_\text{snapshots}$, and $\mathrm{cond}(V)$ characteristics. Measurements via \texttt{scripts/benchmark\_modal\_evolution.py}.

\begin{table}[htbp]
\centering
\caption{Per-mode matrix-exponential evolution benchmarks. Wall time is total per-call (per-block time $\times$ \#blocks) on a single CPU core. Path~(d), the empirically-best choice, is unconditional in v0.31+. All four paths agree to $\sim10^{-12}$ relative.}
\label{tab:modal-expm-benchmark}
\setlength{\tabcolsep}{4pt}
\begin{tabular}{@{}lcccccccccc@{}}
\toprule
Theory & $bs$ & \#blk & $n_\text{mod}$ & $n_\text{snap}$ & (a) eigen & (b) ppm-em & (c) bch-em & (d) Pad\'{e}+mv & (b)/(a) & (d)/(a) \\
\midrule
\texttt{dark\_photon\_plasma} (CDT) & 6 & 5 & 64 & 2 & 19\,ms & 548\,ms & 629\,ms & \textbf{20\,ms} & 28.6$\times$ & \textbf{1.06$\times$} \\
\texttt{torsion\_dark\_photon\_fv} & 6 & 3 & 64 & 2 & 33\,ms & 265\,ms & 283\,ms & \textbf{16\,ms} & 8.0$\times$ & \textbf{0.49$\times$} \\
\texttt{graviton\_torsion} & 10 & 4 & 128 & 50 & 113\,ms & 90{,}627\,ms & 7{,}374\,ms & \textbf{111\,ms} & 803$\times$ & \textbf{0.99$\times$} \\
\texttt{euler\_heisenberg} & 5 & 2 & 96 & 50 & 38\,ms & 17{,}951\,ms & 2{,}270\,ms & \textbf{29\,ms} & 469$\times$ & \textbf{0.75$\times$} \\
\bottomrule
\end{tabular}

\smallskip
\textit{Legend:} (a) eigendecomposition + $V\,\mathrm{diag}(e^{\lambda t})\,V^{-1}$ (original default); (b) per-mode \texttt{expm\_multiply}, one call per (mode,\,snapshot); (c) per-mode \texttt{expm\_multiply} with \texttt{start/stop/num} snapshot scheduling; (d) \texttt{scipy.linalg.expm($M \cdot \mathrm{d}t$)} precomputed once per mode + matvec in the snapshot loop.
\end{table}

\subsubsection{Why path (d) wins on small per-mode blocks}
\label{sec:modal-pade-vs-expm-multiply}

\texttt{scipy.sparse.linalg.expm\_multiply} (paths b, c) carries substantial per-call setup overhead: estimating $\|A\|_1$, choosing scaling parameter $s$, choosing Taylor truncation $m$, allocating Krylov-Taylor workspace. For our per-mode block sizes ($bs \le 15$), this overhead vastly dominates the actual $O(bs^3) \approx 1000$ flops of the exponential. Calling it once per mode (path b) or once per mode with snapshot scheduling (path c) means we pay that overhead $O(n_\text{modes})$ times.

By contrast, \texttt{scipy.linalg.expm} (path d) is a simpler direct routine that computes the matrix exponential of an $(bs \times bs)$ matrix in one $O(bs^3)$ Pad\'{e} call. Once $\exp(M \cdot \mathrm{d}t)$ is precomputed per mode, the snapshot loop becomes pure matvec at $O(bs^2)$ per snapshot per mode, and the snapshot-time iteration is structurally identical to the eigendecomposition's vectorised einsum:
\begin{verbatim}
M_modes = solve(B_modes, A_modes)               # (n_modes, bs, bs), one solve
exp_M_dt = stack(scipy.linalg.expm(M_modes[m] * dt) for m)  # (n_modes, bs, bs)
y_curr = y0_block.copy()
for ti in range(n_snapshots):
    snapshots[ti] = y_curr
    y_curr = einsum("mij,mj->mi", exp_M_dt, y_curr)
\end{verbatim}

This contrasts with \cref{sec:modal-perf-optimizations} Round~1, where the position-dependent path uses \texttt{expm\_multiply} on the full convolution matrix. There the matrix dimension is $n_\text{total} \approx n_\text{slots} \cdot n_\text{modes}$ ($\sim 10^3\text{--}10^4$), at which the per-call setup cost is amortised and \texttt{expm\_multiply}'s sparse-matrix optimization wins. The two paths --- per-mode for constant coefficients (Pad\'{e} + matvec, this section) and full-matrix for position-dependent coefficients (\texttt{expm\_multiply}, \cref{sec:modal-opt-round1}) --- are not in competition; they are tuned to different matrix-size regimes within the same canonical robust-matrix-exponential strategy.

\subsubsection{No threshold gate; eigendecomposition opt-in for Pass 1}
\label{sec:modal-pade-pass1}

Because path~(d) is wall-time-competitive with eigendecomposition \emph{and} mathematically robust for arbitrary $\mathrm{cond}(V)$, the original hybrid plan (\texttt{cond(V) > 1e12} threshold $\Rightarrow$ \texttt{expm\_multiply} fallback) was unnecessary. Path~(d) is unconditional for default Pass~0 evolution.

The eigendecomposition path remains opt-in for the Pass~1 perturbative-Duhamel solver (\cref{sec:perturbative-reduction}), which requires per-eigenvalue closed-form kernels $G(\lambda_i, \lambda_j; t)$ and the $V$, $V^{-1}$ projection. Callers requesting eigendata via \texttt{collect\_eigendata} retain the original codepath, with an added $\mathrm{cond}(V) > 10^{12}$ check that raises \texttt{NotImplementedError} when Pass~0 is unsound (preventing silent corruption of Pass~1 corrections). The follow-up issue \texttt{\#320-Pass1} tracks an augmented-exponential Pass~1 path (Al-Mohy \& Higham 2011 \S5.2: $\exp(t \cdot \begin{psmallmatrix} A & M_\text{src} \\ 0 & A \end{psmallmatrix}) \cdot \begin{psmallmatrix} 0 \\ y^0(0) \end{psmallmatrix}$) that removes the eigendata dependency entirely, leaving \texttt{scipy.linalg.expm} as the sole canonical Pass~0/Pass~1 backend.

\subsubsection{Verification}
\label{sec:modal-pade-verification}

\begin{itemize}
  \item \textbf{Synthetic ill-conditioned unit test} (\texttt{tests/test\_modal\_robust\_evolution.py}): a $4\times 4$ generalized-eigenvalue problem with $\mathrm{cond}(V) \approx 10^{15}$ by construction. Path~(d) gives $\exp(M \cdot t) \cdot y_0$ to $\sim 10^{-12}$ relative; the original $V \cdot \mathrm{diag} \cdot V^{-1}$ path does not.
  \item \textbf{CDT vs.\ FV bit-exact regression}: at five random parameter points spanning the v3 amplify prior, CDT (path d) and FV (well-conditioned, also path d) give $P_\text{max}$ matching to $\Delta / P_\text{max} < 10^{-10}$.
  \item \textbf{No regression on well-conditioned theories}: the existing test suite (\texttt{euler\_heisenberg}, \texttt{graviton\_torsion}, FV scalar/vector examples) passes unchanged --- path (d) wall time is within $1.06\times$ of eigendecomposition on every workload, and bit-exact agreement to $\sim 10^{-12}$.
\end{itemize}

\subsubsection{Strategic implication for Stages B--E}

The cross-formulation cross-check (FV vs.\ CDT) was the empirical signal that triggered this investigation. For more complex PGT theories in the campaign roadmap (Stages~B--E), an FV-style alternative formulation may not exist for cross-checking. Without the unconditional path~(d) fix, those campaigns would silently produce wrong answers in any region of parameter space where the redundant tensor components drive $\mathrm{cond}(V)$ above $10^{13}$. The unconditional fix is therefore a campaign-wide robustness investment, not just a Stage~A bug fix.

\textbf{References}: Higham, N.J.\ (2008). \emph{Functions of Matrices: Theory and Computation}, SIAM, \S2.3 (eigendecomposition abandonment threshold); Higham, N.J.\ (2009). ``The scaling and squaring method for the matrix exponential revisited.'' \emph{SIAM Review} 51(4):747--764; Trefethen, L.N.\ \& Bau, D.\ (1997). \emph{Numerical Linear Algebra}, SIAM, Lecture~33; Al-Mohy, A.H.\ \& Higham, N.J.\ (2011). ``Computing the action of the matrix exponential.'' \emph{SIAM J.\ Sci.\ Comput.} 33(2):488--511; Davies, P.I.\ \& Higham, N.J.\ (2003). ``A Schur--Parlett algorithm for computing matrix functions.'' \emph{SIAM J.\ Matrix Anal.\ Appl.} 25(2):464--485.

\subsubsection{Snapshot output precision}
\label{sec:modal-snapshot-precision}

The modal solver produces two structurally different snapshot output paths, with different precision characteristics.

\paragraph{Uniform step-by-step path (default).}
When the requested snapshot times are equally spaced (\texttt{t\_eval = [0, dt, 2dt, \ldots]}), the solver precomputes $E \equiv \exp(M \cdot \mathrm{d}t)$ once per mode block via \texttt{scipy.linalg.expm}, then applies it iteratively:
\begin{equation}
y_k = E \cdot y_{k-1} = E^k \cdot y_0
\end{equation}
Snapshot $k$ therefore accumulates $k-1$ matrix-multiplication roundoff errors beyond the initial Pad\'{e} backward error. For $k$ steps with block size $bs$, the bound is
\begin{equation}
\|\delta y_k\| \lesssim (k-1) \cdot u \cdot \|E\|^2 \cdot \|y_0\|
\end{equation}
where $u \approx 2.2 \times 10^{-16}$ is IEEE-754 unit roundoff. For the D2.2 campaign configuration ($t_\mathrm{end}=10$, $N=21$ snapshots, $\mathrm{d}t=0.5$), the accumulated error at the final snapshot is ${\sim}19 \times u \approx 4 \times 10^{-15}$ --- negligible for all practical physics measurements.

\paragraph{Non-uniform path.}
When snapshot times are unequally spaced, each snapshot is computed independently from the initial state:
\begin{equation}
y(t_k) = \exp\!\bigl(M \cdot (t_k - t_0)\bigr) \cdot y_0
\end{equation}
Each snapshot is machine-precision at every $t_k$ with no accumulation. The cost is $n_\mathrm{snaps}$ \texttt{expm} calls per mode block, compared to one \texttt{expm} call for the uniform path.

\paragraph{$N=2$ snapshot special case (inference).}
When \texttt{--snapshots=t\_end} is passed (or equivalently \texttt{t\_eval = [0, t\_end]}), the uniform path applies with $\mathrm{d}t = t_\mathrm{end}$: $E = \exp(M \cdot t_\mathrm{end})$ is precomputed once and applied \emph{once}. This is machine-precision at $t_\mathrm{end}$ with zero accumulated error --- identical precision to the non-uniform path at every snapshot, at the cost of a single \texttt{expm} precomputation. This is the recommended configuration for PolyChord/nested-sampling inference runs where only the final state is needed (see \cref{sec:inference-snapshot-guidance}).

\paragraph{Why not switch the uniform path to compute each snapshot from $y_0$?}
For visualization workloads with $N=21$ snapshots and $n_\mathrm{modes}=128$ blocks (e.g.\ \texttt{graviton\_torsion}), switching to direct-from-$y_0$ would require $21 \times 128 = 2688$ \texttt{expm} calls vs.\ $128$ for the current uniform path --- approximately $21\times$ more \texttt{expm} calls. The resulting regression in visualization sweep wall time is not justified by the $4 \times 10^{-15}$ precision improvement. For inference runs where this matters, the $N=2$ configuration already achieves machine precision. The deliberate decision to retain the uniform path is documented in GitHub issue \#331.

\subsection{Performance Optimizations}
\label{sec:modal-perf-optimizations}

The modal solver has undergone two rounds of targeted optimization, each measured with incremental benchmarking. All changes preserve machine-precision accuracy.

\subsubsection{Round 1: Sparse Convolution Matrix}
\label{sec:modal-opt-round1}

\textbf{Problem:} The position-dependent path built a dense $n_\text{total}\times n_\text{total}$ convolution matrix, but for localized coefficients (Gaussian $\Bzero$) most entries are negligibly small.

\textbf{Fix:} Threshold entries below $10^{-14} \times \max|A|$, convert to sparse CSC when density $< 30\%$. \texttt{scipy.sparse.linalg.expm\_multiply} natively supports sparse matrices.

\begin{table}[htbp]
\centering
\caption{Round 1 optimization: sparse convolution matrix.}
\label{tab:modal-opt-round1}
\begin{tabular}{@{}lccccl@{}}
\toprule
Grid $N$ & Density & Dense & Sparse & Speedup & Max diff \\
\midrule
128 & 4.0\% & 5.44\,s & 0.23\,s & $\mathbf{24\times}$ & $7.4\times 10^{-14}$ \\
256 & 2.6\% & 15.85\,s & 0.59\,s & $\mathbf{27\times}$ & $1.2\times 10^{-13}$ \\
512 & 1.4\% & 29.08\,s & 1.23\,s & $\mathbf{24\times}$ & $1.1\times 10^{-12}$ \\
\bottomrule
\end{tabular}
\end{table}

\subsubsection{Round 2: Micro-Optimizations}
\label{sec:modal-opt-round2}

\begin{table}[htbp]
\centering
\caption{Round 2 micro-optimizations.}
\label{tab:modal-opt-round2}
\begin{tabular}{@{}llll@{}}
\toprule
ID & Optimization & Location & Impact \\
\midrule
O1 & Vectorize convolution coupling inner loop & \texttt{\_add\_convolution\_coupling} & pos-dep $N=512$: 1.79$\to$1.36\,s \\
O3 & Batch $S_{cc}$ inversion (eliminate per-mode loop) & \texttt{\_build\_constraint\_eliminated\_matrices} & constraint path $1.7\times$ faster \\
O4 & Batch vel\_coupling solve & \texttt{\_build\_constraint\_eliminated\_matrices} & (part of constraint $1.7\times$) \\
O5 & Analytical rfft shape (eliminate FFT-of-zeros probe) & \texttt{\_fft\_slots} & removes wasted computation \\
O6 & Pre-allocate \texttt{y\_hat\_t} buffer outside timestep loop & \texttt{\_evolve\_per\_mode} & fewer allocations \\
O7 & Return Fourier snapshots to avoid re-FFT & \texttt{\_evolve\_per\_mode} + \texttt{solve\_modal} & constraint path $1.7\times$ faster \\
\bottomrule
\end{tabular}
\end{table}

\textbf{Rejected:} O2 (vectorize snapshot evaluation across all output times) --- the batched exp+einsum created larger intermediate arrays with more memory pressure, making it $1.3\times$ slower than the lightweight per-timestep loop.

\subsubsection{Combined Result}
\label{sec:modal-opt-combined}

\begin{table}[htbp]
\centering
\caption{Combined optimization results across all modal solver paths.}
\label{tab:modal-opt-combined}
\begin{tabular}{@{}llll@{}}
\toprule
Path & Before optimization & After & Total speedup \\
\midrule
Constant-coeff (eigendecomp, $N=512$) & 0.026\,s & 0.020\,s & $\mathbf{1.3\times}$ \\
Position-dep (\texttt{expm\_multiply}, $N=512$) & ${\sim}29$\,s & 1.4\,s & $\mathbf{{\sim}21\times}$ \\
Constraint (Proca 3D, $16^3$) & 0.136\,s & 0.078\,s & $\mathbf{1.7\times}$ \\
\bottomrule
\end{tabular}
\end{table}

\subsection{References}
\label{sec:modal-references}

\begin{itemize}
  \item Moler, C.\ \& Van Loan, C.\ (2003). ``Nineteen dubious ways to compute the exponential of a matrix, twenty-five years later.'' \emph{SIAM Review}, 45(1):3--49.
  \item Golub, G.H.\ \& Van Loan, C.F.\ (1996). \emph{Matrix Computations}, 3rd ed.\ Johns Hopkins University Press.
  \item Hairer, E., Lubich, C.\ \& Wanner, G.\ (2006). \emph{Geometric Numerical Integration}, Springer, \S4.
  \item Burns, K.J.\ et al.\ (2020). ``Dedalus: A flexible framework for numerical simulations with spectral methods.'' \emph{Phys.\ Rev.\ Research}, 2:023068.
  \item Raffelt, G.\ \& Stodolsky, L.\ (1988). ``Mixing of photons with low-mass particles in magnetic fields.'' \emph{Phys.\ Rev.\ D}, 37:1237.
  \item Dormand, J.R.\ \& Prince, P.J.\ (1980). ``A family of embedded Runge-Kutta formulae.'' \emph{J.\ Comp.\ Appl.\ Math.}, 6.
  \item Hindmarsh, A.C.\ et al.~\cite{hindmarsh2005sundials}. SUNDIALS: Suite of nonlinear and differential/algebraic equation solvers.
  \item Al-Mohy, A.H.\ \& Higham, N.J.\ (2011). ``Computing the action of the matrix exponential.'' \emph{SIAM J.\ Sci.\ Comput.}, 33(2):488--511.
  \item Trefethen, L.N.\ \& Embree, M.\ (2005). \emph{Spectra and Pseudospectra}. Princeton University Press.
  \item Van Loan, C.\ (1978). ``Computing integrals involving the matrix exponential.'' \emph{IEEE Trans.\ Automatic Control}, 23(3):395--404.
  \item Higham, N.J.\ (2008). \emph{Functions of Matrices: Theory and Computation}. SIAM. \S2.3, Ch.\ 12.
  \item Higham, N.J.\ (2009). ``The scaling and squaring method for the matrix exponential revisited.'' \emph{SIAM Review}, 51(4):747--764.
  \item Trefethen, L.N.\ \& Bau, D.\ (1997). \emph{Numerical Linear Algebra}. SIAM. Lecture~33.
  \item Davies, P.I.\ \& Higham, N.J.\ (2003). ``A Schur-Parlett algorithm for computing matrix functions.'' \emph{SIAM J.\ Matrix Anal.\ Appl.}, 25(2):464--485.
  \item Yoshida, H.~\cite{yoshida1990symplectic}. Construction of higher-order symplectic integrators.
\end{itemize}
